\section{Introduction}

Interpretability is increasingly used not only to explain models but also to
change what they do. An estimate of a model's internal state may guide an
activation edit, decide which components to disable, determine whether a tool
call is allowed, or rank requests for a limited audit budget. Activation
steering and model-based control make this use explicit
\cite{turner2023activation,zou2023representation,nguyen2025feedback,skifstad2026local}.
Once an estimate helps choose an action, prediction quality is only part of the
evaluation.

We call a method that supplies such an estimate an \emph{observer}. We call the
rule that turns the estimate into an action a \emph{controller}. An
\emph{intervention} is the edit or decision that the controller is allowed to
make. This paper asks a simple question: is the observer good enough for the
action the controller must choose?

Three examples show why prediction and action must be tested separately.

\begin{itemize}
  \item Two requests may both violate policy, but an unauthorized delete can
  cause more harm than an unauthorized read. A classifier can label both
  correctly and still use a limited block or review budget poorly.

  \item An observer may correctly predict the average effect of disabling a
  set of model components. The same edit may do nothing on some prompts and
  overshoot on others, making it a poor choice for a prompt-level target.

  \item A classifier built from an internal model state may appear informative
  while repeating a signal already present in the prompt or output. It then
  adds no useful information at the point where the decision is made.
\end{itemize}

We introduce \textbf{ObserverBench}\footnote{Public workbench and submission
page: \url{https://kwisatzh.github.io/observerbench/} (release forthcoming).},
a set of fixed, reusable tasks for making these distinctions measurable. Each
task specifies the model, target,
information available to the observer, allowed interventions, action rule,
measurement budget, unseen test cases, and loss. The same rule is applied to
every submitted observer. The evaluator reports both how accurately the
observer predicts and how well the resulting actions perform.

ObserverBench addresses a different stage from mechanistic tomography. The
mechanistic tomography framework develops measurement designs for recovering
internal quantities and intervention effects; ObserverBench starts with a
proposed estimate and asks whether it is adequate for a fixed intervention,
control, or safety use \cite{erramilli2026tomography}.

An ObserverBench result is defined by five choices:

\begin{itemize}
  \item \textbf{What does the observer estimate?} A model feature, the average
  effect of an edit, and the cost of an action on one prompt are different
  targets.

  \item \textbf{What information may it use?} Prompt text, internal model
  states, and output scores provide different information and have different
  costs.

  \item \textbf{What actions are allowed?} An error may not matter if it lies
  outside the part of model behavior that the allowed edits can affect.

  \item \textbf{How is the action chosen and scored?} The same estimate can be
  useful under one decision rule or loss and poor under another.

  \item \textbf{Where will it be used?} Prompt families, attack frequency, and
  intervention budgets can change which observer performs best.
\end{itemize}

These choices make the result task-relative. An observer can differ from the
target in ways that never affect an allowed action. A method can predict an
average edit more accurately without choosing a better prompt-level edit. A
safety score can rank policy labels correctly while ignoring the different
consequences of those labels. System boundaries matter as well: recent cyber
evaluations show that tools, permissions, external state, and side channels
can affect what a prompt-level score means
\cite{openai2026securityincident,huggingface2026securityincident}.

\paragraph{Three evaluation modes.}
The current release applies this idea in three settings:

\begin{itemize}
  \item \textbf{Control.} The task fixes the model, starting states,
  controller, interventions, and loss. A comparison changes the state
  estimator or edit direction while keeping the rest of the loop fixed.

  \item \textbf{Effect prediction.} The task provides fixed measurements and
  unseen interventions. A submission predicts their effects; the evaluator
  then applies a fixed action rule and reports prediction error and action
  loss separately.

  \item \textbf{Safety triage.} A submission provides one risk score per
  request or trajectory. The evaluator turns those scores into allow, block,
  escalate, or audit decisions under a fixed budget and measures both missed
  harm and unnecessary intervention.
\end{itemize}

\paragraph{Contributions.}

\begin{itemize}[itemsep=0.2em,topsep=0.3em]
  \item \textbf{A workbench researchers can extend.} A researcher submits an
  observer as code or a prediction table. A fixed task reports both prediction
  quality and the cost of the resulting actions.

  \item \textbf{A formal adequacy test for control.} For linear readouts with
  an offset, two
  observers need to agree only at the starting point and along directions the
  allowed edits can reach. A fixed-step formula separates errors in the
  starting value, response to an edit, and changes the observer cannot see.

  \item \textbf{Evidence that prediction and action can disagree.} On
  GPT-2-small and Qwen2.5-7B, pairwise terms improve unseen effect prediction
  without always improving the fixed action, while direct-loss observers
  lower downstream loss.

  \item \textbf{Safety scores tested as decisions.} The safety tasks compare
  observers under fixed block, review, or audit budgets. Correct label ranking
  can still spend that budget poorly, and realistic attack rates can separate
  methods that a balanced test makes appear tied.
\end{itemize}
